% =============================================================================
% 学术论文模板 - 支持中英文
% =============================================================================

% 文档类选择
\documentclass[12pt, a4paper]{article}

% =============================================================================
% 宏包加载
% =============================================================================

% 中文支持 (使用 XeLaTeX 编译)
\usepackage{xeCJK}
\setCJKmainfont{Songti SC}      % 宋体
\setCJKsansfont{Heiti SC}       % 黑体
\setCJKmonofont{PingFang SC}    % 苹方（等宽）

% 字体与编码
\usepackage[T1]{fontenc}
\usepackage{mathptmx}         % Times 字体
\usepackage{helvet}

% 语言支持
\usepackage[english]{babel}
\usepackage{csquotes}         % 引用格式

% 页面布局
\usepackage[a4paper, margin=2.5cm]{geometry}
\usepackage{parskip}          % 段落间距

% 图表支持
\usepackage{graphicx}
\usepackage{float}
\usepackage{caption}
\usepackage{subcaption}

% 表格增强
\usepackage{booktabs}         % 专业表格线
\usepackage{tabularx}
\usepackage{multirow}

% 数学公式
\usepackage{amsmath, amssymb, amsfonts}
\usepackage{mathtools}
\usepackage{physics}          % 物理符号

% 单位
\usepackage{siunitx}

% 代码高亮
\usepackage{listings}
\usepackage{xcolor}

% 算法伪代码
\usepackage[ruled,vlined,linesnumbered]{algorithm2e}

% 参考文献设置
\usepackage[backend=biber,style=ieee]{biblatex}
\addbibresource{references.bib}

% 超链接（必须在 cleveref 之前）
\usepackage{hyperref}
\hypersetup{
    colorlinks=true,
    linkcolor=blue,
    citecolor=blue,
    urlcolor=blue,
}

% 交叉引用（必须在 hyperref 之后）
\usepackage{cleveref}
\crefname{section}{Section}{Sections}
\crefname{figure}{Figure}{Figures}
\crefname{table}{Table}{Tables}

% 术语表
\usepackage[acronym]{glossaries}
\makeglossaries

% =============================================================================
% 代码高亮设置
% =============================================================================
\definecolor{codegreen}{rgb}{0,0.6,0}
\definecolor{codegray}{rgb}{0.5,0.5,0.5}
\definecolor{codepurple}{rgb}{0.58,0,0.82}
\definecolor{backcolour}{rgb}{0.95,0.95,0.92}

\lstdefinestyle{mystyle}{
    backgroundcolor=\color{backcolour},
    commentstyle=\color{codegreen},
    keywordstyle=\color{magenta},
    numberstyle=\tiny\color{codegray},
    stringstyle=\color{codepurple},
    basicstyle=\ttfamily\footnotesize,
    breakatwhitespace=false,
    breaklines=true,
    captionpos=b,
    keepspaces=true,
    numbers=left,
    numbersep=5pt,
    showspaces=false,
    showstringspaces=false,
    showtabs=false,
    tabsize=2
}
\lstset{style=mystyle}

% =============================================================================
% 标题信息
% =============================================================================
\title{学术论文标题\\
\large 学术论文副标题}
\author{作者姓名}

\date{\today}

% =============================================================================
% 文档开始
% =============================================================================
\begin{document}

% 标题页
\maketitle

% 摘要
\begin{abstract}
这是一篇学术论文的摘要。摘要应该简明扼要地介绍研究背景、目的、方法、结果和结论。本模板支持中英文混排，包含图表、公式、代码和算法等常用功能。

关键词：LaTeX、学术论文、模板、中文支持
\end{abstract}

% 目录
\newpage
\tableofcontents
\newpage

% =============================================================================
% 正文
% =============================================================================

\section{引言}

\subsection{研究背景}
这是引言部分。这里介绍研究的背景和意义~\cite{example2024}。

\subsection{研究目的}
本研究的主要目的是提供一个完整的学术论文模板，包含以下功能：
\begin{itemize}
    \item 中英文混排支持
    \item 图表自动编号和引用
    \item 数学公式支持
    \item 代码高亮
    \item 算法伪代码
    \item 参考文献管理
\end{itemize}

\section{相关技术}
\label{sec:related}

相关技术介绍...参见 \cref{fig:example}。

\section{方法}

\subsection{数学公式}
\label{sec:math}

本节展示数学公式的使用方法。

\subsubsection{行内公式}
爱因斯坦质能方程 $E=mc^2$ 是物理学中最著名的公式之一。

\subsubsection{行间公式}
薛定谔方程：
\begin{equation}
    i\hbar\frac{\partial}{\partial t}\Psi(\mathbf{r},t) = \hat{H}\Psi(\mathbf{r},t)
    \label{eq:schrodinger}
\end{equation}

\subsubsection{多行公式}
\begin{align}
    f(x) &= x^2 + 2x + 1 \notag \\
         &= (x+1)^2 \label{eq:factorization}
\end{align}

\subsection{算法伪代码}
\label{sec:algorithm}

\begin{algorithm}
    \SetAlgoLined
    \KwIn{输入参数}
    \KwOut{输出结果}
    初始化变量\;
    \While{条件满足}{
        执行操作\;
        \If{特殊情况}{
            处理特殊情况\;
        }
    }
    返回结果\;
    \caption{示例算法}
    \label{algo:example}
\end{algorithm}

\subsection{代码展示}
\label{sec:code}

\begin{lstlisting}[language=Python, caption=Python 代码示例]
def hello_world():
    """打印 Hello World"""
    print("Hello, World!")
    return 0

if __name__ == "__main__":
    hello_world()
\end{lstlisting}

\section{实验结果}

\subsection{表格}
\label{sec:table}

\begin{table}[H]
    \centering
    \caption{实验结果对比表}
    \label{tab:results}
    \begin{tabular}{lccc}
        \toprule
        方法 & 准确率 & 召回率 & F1 分数 \\
        \midrule
        方法 A & 85.3 & 82.1 & 83.7 \\
        方法 B & 87.6 & 84.5 & 86.0 \\
        \textbf{本文方法} & \textbf{90.2} & \textbf{88.7} & \textbf{89.4} \\
        \bottomrule
    \end{tabular}
\end{table}

结果见 \cref{tab:results}，本文方法表现最佳。

\subsection{图片}
\label{sec:figure}

\begin{figure}[H]
    \centering
    % 示例图片 - 替换为实际图片路径
    % \includegraphics[width=0.8\textwidth]{figures/example.png}
    \rule{0.8\textwidth}{0.4\textwidth} % 占位符
    \caption{示例图片}
    \label{fig:example}
\end{figure}

如 \cref{fig:example} 所示...

\subsection{多图排版}

\begin{figure}[H]
    \centering
    \begin{subfigure}{0.45\textwidth}
        \centering
        \rule{0.9\textwidth}{0.3\textwidth}
        \caption{子图 (a)}
    \end{subfigure}
    \hfill
    \begin{subfigure}{0.45\textwidth}
        \centering
        \rule{0.9\textwidth}{0.3\textwidth}
        \caption{子图 (b)}
    \end{subfigure}
    \caption{多图排版示例}
\end{figure}

\section{讨论}

对比 \cref{eq:schrodinger} 和 \cref{eq:factorization}...

算法流程见 \cref{algo:example}，代码实现见 \cref{sec:code}。

\section{结论}

本文提供了一个完整的学术论文模板...

% =============================================================================
% 参考文献
% =============================================================================
\newpage
\printbibliography

% =============================================================================
% 附录
% =============================================================================
\appendix
\section{附录材料}

附录内容...

\end{document}
